% !TEX root = ../main.tex
\chapter{Performance Evaluation}
\label{ch:evaluation}

% =====================================================================
% ★ 內容權威來源：docs/thesis/evaluation-protocol-v1.md ★
% 那份文件是完整的協定，本章是它的正文化。改動請先改那份，再同步過來。
% 舊的 thesis-outline.md §7.2–7.4 已作廢，不要回頭參照。
%
% 【2026-08-23 動筆註】§5.1／§5.2／§5.3／§5.5／§5.10 已寫成正文。體例沿用第 4 章：
% 每一節末尾一個 "Status:" 段落，主張要嘛 implementation-backed（指名檔案／實測數字），
% 要嘛 normative（明說尚未實證並揭露偏誤方向）。不允許沒有標記的中間態——第 3 章
% 骨架漂移的六處全部出在沒有標記的段落上。
% 尚未動筆的 §5.4／§5.6／§5.7／§5.8／§5.9 保留原本的 % TODO 寫作規格。
% =====================================================================

\section{Research Questions}
\label{sec:eval:rq}

The claim under test is not that a language model can fill in a form. It is that, for a
bounded class of recurring template-instantiation tasks, the model can be removed from the
recurring execution path altogether without losing correctness, traceability, or fail-loud
behaviour. Four research questions decompose that claim; a fifth asks which components of
the derivation procedure carry it.

\begin{description}
  \item[RQ1 --- Correctness.] Is a compiled specification's field-level semantic accuracy at
    least as good as a per-instance agent's on the same instances, measured by the same
    deterministic evaluator?
  \item[RQ2 --- Consistency.] How much does repeated execution of the same task vary? The
    question is asymmetric by construction, and the asymmetry is the point: for a
    per-instance agent the variance is present in every one of the $N$ executions, whereas
    for the compiled method execution is deterministic and \emph{all} of the variance sits in
    the single derivation. RQ2 therefore has two halves, and reporting only the first would
    be a misstatement: run-to-run variation of the baselines, and derivation-to-derivation
    variation of the compiled specifications.
  \item[RQ3 --- Cost structure.] What is the shape of cost against instance count? A one-off
    derivation cost plus a near-zero marginal execution cost crosses a per-instance cost at
    some $N$; the question is where that crossing lies, and how sensitive it is to derivation
    retries.
  \item[RQ4 --- Behaviour under template drift.] When the template changes underneath a
    frozen artifact, does the method fail loudly and locate the change, and do the baselines
    fail silently? This is the direct test of the design principle that governs the whole
    system.
  \item[RQ5 --- Component contribution.] Which of the structural gates, the two oracles, and
    the repair loop are load-bearing, and which are redundant with each other?
\end{description}

An earlier version of this plan carried a sixth question --- whether the method generalizes
beyond drawings to office documents. It has been removed rather than answered, because DOCX
and XLSX are now the primary evaluation carriers rather than a proof of concept, and a
question whose answer is a premise of the study design is not a research question.

\paragraph{A revision to the emphasis, and the measurement that forced it.}
The original framing put RQ1 first and treated RQ2 and RQ3 as secondary. A pilot measurement
on the DOCX dataset changed that ordering. Running a per-instance agent three times over one
calibration instance produced semantically correct output every time --- the cross-file
lookup, the rounding, the gram conversion, and the extent decision were all right --- and
three different files: three distinct output hashes in three runs, and cell-level agreement
of 119 of 172 cells, or 69.2\%. All 53 divergent cells were formatting decisions retaken from
scratch on each run, a thousands separator present in one run and absent in the next, and
none of them was a semantic error.

The consequence is stated here rather than in the results, because it is a claim about what
this thesis argues: the primary contribution is not that the compiled method is \emph{more
accurate} than a competent agent on tasks a competent agent can do. It is that accuracy is
comparable while consistency and cost structure differ structurally. Two tempting responses
to that pilot were considered and rejected. Raising task difficulty until the baseline fails
would tune the benchmark towards making the comparator fail; substituting a weaker model
would be choosing the opponent. Both are validity threats, and both are recorded as rejected
in Section~\ref{sec:eval:threats}.

\paragraph{Status: implementation-backed as a measurement, normative as a research plan.}
The pilot above was executed and its artifacts are retained. The five questions are a plan;
the formal experiment has not been run, and no figure in this chapter should be read as an
answer to any of them.

\section{Study Design}
\label{sec:eval:design}

\subsection{Why Not a Few-Shot Held-Out Split}
\label{sec:eval:whynotfewshot}

The conventional evaluation for a language-model system is a few-shot held-out split: the
model sees $K$ demonstrations and is scored on unseen tasks. That protocol borrows its
meaning from in-context learning, and this method is not that shape.

The template is an \emph{input}, not the object of generalization. The system is given the
template on purpose; so is every baseline, on every one of its executions. Seeing the
template is therefore not leakage but a premise, and it is symmetric across arms. This makes
the standard split either vacuous or unfair depending on how it is drawn. If the held-out
instances share a template with the demonstrations, nothing was held out that matters: the
compiled specification and the agent's demonstrations both describe that template. If the
held-out instances come from a different template, the compiled method must derive again ---
the demonstrations do it no good at all --- and the comparison measures the cost of a
derivation rather than generalization from examples.

What the split is genuinely needed for is narrower, and is retained: preventing the situation
in which one arm's prompts, gates, or intermediate representation were tuned on the very data
it is scored on while the baselines were not. The instrument for that is separation at the
\emph{template} level, and Section~\ref{sec:eval:split} reports how far that separation
actually holds in this study.

\paragraph{Status: normative, and deliberately so.} This subsection argues a design choice.
It is included because the objection it answers is the first one a reader of a
language-model-systems paper will raise, and because declining a standard protocol without
stating why reads as avoiding it.

\subsection{Two-Level Split}
\label{sec:eval:split}

The split has two levels, and only the first of them currently holds in full.

\paragraph{Instance level.}
Every evaluation template carries ten instances: $K=3$ calibration instances, shared by all
three arms and the only instances any arm may see during construction; two development
instances used for debugging; and $M=5$ sealed evaluation instances that no arm sees at any
stage. Two of the five evaluation instances in each dataset are cases whose \emph{correct}
behaviour is a refusal, so that an arm which always produces a file cannot score well merely
by producing files.

\paragraph{Template level.}
The evaluation templates are two, with a third planned:

\begin{itemize}
  \item \textbf{S4} (DOCX): a Commercial Invoice and a Packing List produced together from a
    product master file and a shipment detail file. Its principal difficulty is cross-file
    computation --- unit price from one workbook, quantity from another --- and it is the only
    DOCX evaluation dataset, so every DOCX claim in this thesis rests on it. Thirty-four
    applicable rules.
  \item \textbf{X1} (XLSX): the Directorate General of Highways monthly site-audit report
    chain, in which $N$ real 13020A audit sheets are aggregated into one 13020D monthly
    report. Its principal difficulties are a cross-file collection of source artifacts and a
    layout in which one logical record occupies two physical rows. Twenty-seven applicable
    rules.
  \item \textbf{S5} (XLSX, not yet built): a credit-programme completion audit, chosen for
    its density of negative obligations.
\end{itemize}

Two further items must be named, because their absence changes how the split reads.
\textbf{S3} is a synthetic pilot dataset that was built to protocol and then rejected on the
grounds that its layout and content were not those of a document that really circulates; it
is retained only as a pilot for the evaluator and the range locator, is not in the main
results table, and is not offered as evidence of external validity. \textbf{S1} and
\textbf{S2} appear in the protocol document as a proposed assignment of two development
templates and \emph{were never built}.

\paragraph{The gap this leaves, and the direction of its bias.}
Because S1 and S2 do not exist, development did not happen on templates disjoint from the
evaluation set. It happened on the DWG case study, on the rejected pilot, and --- decisively
--- on the calibration instances of S4 and X1 themselves. Three intermediate-representation
primitives were added specifically so that the compiled arm could produce X1 at all: artifact
collections, stacked records with a slot origin, and equality merge between resolved row sets.
X1 is an evaluation template. The template-level separation the protocol describes is
therefore, at present, a normative claim rather than a property of this study, and its bias
runs \emph{in favour of the proposed method}.

Following the third of the three admissible responses to such a gap, the claim is kept and
labelled rather than quietly weakened, and the bias is disclosed again in
Section~\ref{sec:eval:threats}. Two measurements bound its size without cancelling it. First,
each primitive was added only after it had been observed to fail on a real file, and its
collateral effect on the already-stored specifications was swept: the slot-alignment gate
introduced with the third primitive was replayed against all thirty-four stored
specifications and matched exactly the two it was written for, with no false positive on the
five specifications of the other dataset. Second, value-level verification of the resulting
specification was performed on development instances the derivation never saw: the fourth X1
specification reproduced 255 of 255 scored values across the three calibration and two
development instances using the scoring policy's own comparator, and a deliberately corrupted
variant of the same specification scored 240. Neither measurement makes the development
location correct; they bound how far it could have been exploited.

\paragraph{A second construct-validity item, recorded here and again in
Section~\ref{sec:eval:threats}.}
X1's rules require projects to be reported in ascending project-number order. The real 13020D
form prescribes no ordering; the ordering is a convention this dataset introduced so that the
extent has a checkable manifestation. The mechanism that implements it is built and tested,
but the question of \emph{who decided the order} belongs to dataset construct validity rather
than to the method.

\paragraph{Status: implementation-backed at the instance level, normative at the template
level, with the bias direction stated.}

\subsection{Dataset Construction}
\label{sec:eval:dataset}

% TODO 每個案例由人工撰寫的 generator 產生三樣：input.* / expected.* / ground\_truth.json
%      期望值必須從規則文件算出來，不得由受測系統產生——否則比對是循環論證。
% TODO 案例之間明細列數必須不同（3/7/15/40），且至少一個案例帶污染列。
%      否則外延規則根本沒被測到，一份寫死 D5:D19 的規格會全過。（ADR 0002 §4.1）
%
% TODO 標註規則已定案（2026-08-15），權威是 docs/thesis/annotation-manual-v1.md
%      ＋ 兩份 JSON Schema。這一節要寫的是它的兩層結構與為什麼是兩層：
%      每模板一份判準（manifest）、每實例一份值（ground\_truth）。
%      理由不是美觀而是一致性——24 個實例各自重複一次判準，任何一次抄漏都會讓
%      兩個實例用不同標準判分，而那種不一致在主結果表上看起來會像方法的差異。
% TODO 五種必要能力（協定原本列四種，標註定案時補上第五種）：
%      數值容差與格式、negative obligation、範圍外延（逐列列出而非只給起訖）、
%      semantic/presentation 雙層、以及 untouched（沒有任何欄位宣告的位置被動到）。
%      第五種的理由要寫：四種全滿足卻把模板其他地方清空的輸出仍可拿滿分。
% TODO 資料集約束再加一條：每個 aggregate 欄位必須配一個獨立的交叉驗證欄位，
%      理由見 §5.5——外延診斷只有 Arm C 宣告得出來，交叉驗證是唯一三個 arm 對等的偵測手段。

\section{Baselines and Fairness}
\label{sec:eval:baselines}

Three arms are compared. They differ only in what persists between instances.

\begin{description}
  \item[Arm A --- per-instance agent.] For each instance, an agent is given the template, the
    rule document, and the source data, and produces the output document. Nothing persists:
    the next instance starts from the same prompt with no artifact carried over. This is the
    status quo the thesis argues against, and it is deliberately a strong version of it ---
    a current frontier model with the same structured inspection tools the proposed method
    uses.
  \item[Arm B --- agent-built frozen skill.] The agent is first allowed to generalize from
    the calibration instances into a reusable artifact of its own design --- a script, a
    checklist, a set of instructions --- which is then frozen. Every evaluation instance is
    executed by an agent that may read that artifact and may not modify it.
  \item[Arm C --- the proposed method.] One derivation produces a specification; every
    instance is executed by the deterministic compiler with no model call at all.
\end{description}

\paragraph{Why three arms and not two.}
Arm B exists to answer the strongest available objection to the result. If only A and C were
compared, any advantage of C could be attributed to the mere existence of a persistent
artifact rather than to anything about how that artifact is produced and verified. Arm B
holds persistence fixed and varies only the production procedure: an agent generalizing
freely, versus a derivation constrained by an intermediate representation, structural gates,
and two oracles. The question Arm B answers is therefore the one that actually distinguishes
the contribution --- is the verified derivation better than an agent's own attempt at the
same idea?

\paragraph{Fairness has two layers, and the second is the one usually omitted.}
The first layer is \emph{tooling parity}: the same structured dump utilities, the same model,
the same retry ceiling, the same wall-clock ceiling. The second is \emph{material parity}:
Arms A and B receive an input list that is byte-identical to the one Arm C's derivation
receives, because an advantage produced by handing one arm a cleaner or richer set of inputs
is not an advantage of the method. The frozen input manifest is truth-free by construction and
its envelope is hashable, so what each arm was given is recorded rather than asserted.

Material parity has a specific consequence that must be stated because it works against the
proposed method. The rule material given to any arm may not contain target coordinates ---
no cell addresses, no bookmarks, no content-control tags, no drawing handles. Arm C's
derivation therefore has to locate its own write targets, and cannot be given them; the rule
table is not permitted to become an answer key for anyone.

\paragraph{Arm B's freeze must be enforced by the filesystem.}
An instruction not to modify the skill is not a control. The frozen artifact is hashed and
mounted read-only for the evaluation phase; an attempt to write to it is recorded as a
protocol violation and reported rather than silently tolerated. This matters because the
failure it prevents --- an agent quietly improving its own skill between instances --- would
convert Arm B into a per-instance arm while still being reported as a frozen one.

\paragraph{Status: contract implementation-backed, live execution not yet built.}
The three-arm input and evidence contract is frozen and tested: the truth-free input
manifest, the hashable envelope, the Arm B freeze, namespace isolation for Arm C, retention
of raw attempt evidence, and fake-capture negative controls are all implemented. The live
adapters that actually drive Arms A and B, and the formal wrapper around Arm C, are not.
The pilot reported in Section~\ref{sec:eval:rq} was run outside that harness and is labelled
as a pilot for exactly that reason; it must not be read as an Arm A result.

\section{Metrics}
\label{sec:eval:metrics}

% TODO 正確性：semantic 欄位級正確率（主指標）、案例級全對率（含 negative obligation）、
%              沉默失敗率、範圍外延正確率（診斷用）
%      presentation 分開報，不混算。不同載體不可以合併成單一正確率。
% TODO 可靠性：A/B 的 run-to-run 一致率與全 R 次皆正確的案例比例；
%              C 的執行期恆等（雜湊實證）+ 推導期變異度
% TODO 成本：derive 的 token/金額/wall-clock（一次性）、每實例成本、break-even N
% TODO 覆蓋率：以確定性切段後的規則文件為分母（現況見 §4.6）

\subsection{Asymmetric Rerun Budget}
\label{sec:eval:reruns}

% TODO ★ 要主動辯護這個設計 ★
%      A/B 各 R=5；C 主表 R=1 + 指定子集 R=5 以輸出雜湊實證恆等。
%      理由：C 的執行期是零 LLM 確定性的，跑五次必然逐位元相同，那不是實驗是複製檔案。
%
% TODO ★ 但 C 的變異度沒有消失，只是換了位置：它在推導期 ★
%      量法：同一模板獨立 derive D=3 次 -> 3 份規格 -> 每份跑全部評測實例，
%      報告三份規格之間的正確率／覆蓋率分佈與 spec diff 規模。
%      把 C 報成「變異度 0」而不交代這件事，會被審稿人一句話擊破。

\section{The Evaluator}
\label{sec:eval:evaluator}

All three arms are scored by one deterministic program. It parses the produced artifact,
extracts the values at the locations the dataset's annotation manifest names, compares them
against that instance's ground truth, and emits per-check results. It never asks a language
model whether an output looks right, for the reason given in Section~\ref{sec:der:oracles}:
a judge that can be wrong in the same direction as the system under test cannot bound the
system's error.

\subsection{Separate Channels, Never One Average}
\label{sec:eval:channels}

Scores are partitioned into channels, each with its own denominator: \emph{semantic} (the
primary metric), \emph{extent} (which source rows were included), \emph{refusal} (a case
whose correct behaviour is to refuse), \emph{structure} (the artifact opens and has the
expected shape), \emph{presentation} (formatting, reported separately from meaning),
\emph{untouched} (regions no field declares were not modified), and \emph{integrity} (the
package's critical parts survived). A check whose prefix is not one of the declared channels
fails closed rather than being scored as absent. The policy is versioned, and every result
record carries the evaluator version, the policy version, the source hash of the scoring
code, and the repository commit.

The partition is not presentational. The one measurement that settles it: while the ``this
slot must remain blank'' obligations were counted inside the semantic denominator, an arm
that submitted the blank template unchanged --- writing nothing whatsoever --- scored 66.7\%
semantic correctness on one instance. After the obligations were moved into their own
channel the same output scores 0\%. Every individual check had been correct; what was wrong
was which total they were added into.

\subsection{Two Self-Validations, Neither Sufficient Alone}
\label{sec:eval:selfvalidation}

\paragraph{Forward.} The evaluator must score the generator's own expected artifacts at
100\%. This catches an evaluator that reads the wrong location or compares with the wrong
tolerance.

\paragraph{Reverse.} Controlled corruptions are applied to the expected artifact and the
evaluator must catch every one. Without this, an evaluator that always answers ``pass'' would
give every arm a perfect score and the forward test would not notice. The current mutation
coverage is 6 of 6 on the pilot dataset, 12 of 12 on S4, and 13 of 13 on X1; X1 is the first
dataset in which \emph{all seven} channels have at least one mutation that exercises them,
which closed two channels that had previously been structurally vacuous.

\paragraph{A mutation test is not a validity proof, and the gap is specific.}
The expected artifact and the ground truth are produced by one generator. If the generator
computes a value wrongly, both are wrong together and the evaluator faithfully reports 100\%.
Mutation testing establishes that the evaluator detects corruption; it says nothing about
whether the generator computed correctly. The two do not overlap. The present detection
mechanism for that residual is a per-dataset spot audit --- at least three fields per dataset
recomputed by hand from the rule document, including one aggregate and one negative
obligation --- which is probabilistic rather than a guarantee, and is reported as such.

\subsection{A Negative Control for Every Zero}
\label{sec:eval:negativecontrols}

The evaluator and the probes around it are themselves software, and the failure mode observed
repeatedly during construction is that a broken probe reports a serene zero: zero mismatches,
zero problems, full marks. Measured instances during this work include a probe that used the
wrong key name, one that omitted a worksheet prefix from every reference so that nothing
matched anything, one whose trial compilation failed so that the comparison returned an empty
list, and one that compared string forms where the annotation declared exact decimal
comparison, so that twenty \emph{correct} cells were reported as wrong. Two of those report
``looks broken'' and the rest report ``looks fine''; the second kind is the dangerous one, and
it is indistinguishable from success by inspection.

The rule adopted is therefore mechanical rather than advisory. Any reported zero or perfect
score must be accompanied by a deliberate control that makes the figure move, and the control
must distinguish doing the work from not doing it --- a control whose expected value also
happens to be the result of skipping the mechanism measures nothing and never turns red. That
last clause is not hypothetical: an ordering test was found whose expected row order equalled
the arrival order, so disabling the sort entirely left it green.

\subsection{The Extent Diagnostic Is Not Symmetric Across Arms}
\label{sec:eval:extentasymmetry}

Only Arm C declares which source rows it included; Arms A and B produce a document and no
provenance. Scoring the extent channel from an arm's own declaration would repeat exactly the
mistake that the coverage denominator made in Section~\ref{sec:der:coverage} --- a figure
self-reported by the system being scored is not a cross-arm metric.

The cross-arm instrument is therefore a set of \emph{cross-checks}: identities between two
independently scored fields of the same instance, which any arm's output either satisfies or
does not. Alongside them, each dataset's manifest declares, for every wrong-extent hypothesis,
the values that hypothesis would produce, so a wrong extent can be \emph{named} from the
output alone rather than inferred. The effect is measurable: an output that occupies one extra
detail slot while leaving all six totals byte-identical passes all 112 semantic checks on one
instance and is still identified as a specific wrong-extent hypothesis. Arm C's own resolved
anchors are reported as supplementary evidence and are explicitly marked as available for one
arm only; they are never placed in a column alongside A and B.

\paragraph{Status: implementation-backed.}
The shared evaluator lints, self-tests, and scores all three datasets; the scoring policy is
frozen at version 1.1.0; the mutation suites and their per-channel negative controls are in
the repository and run in continuous integration. What is not established is the generator's
own arithmetic, as stated above.

\section{Template Drift}
\label{sec:eval:drift}

% TODO 漂移是獨立的一條軸，dev/eval 切分全是同模板，吃不掉它。
%      一個變體只動一個變因（挪動／改欄位名／換 block／增列擇一），否則失敗時分不出成因。
%      觀察：A/B 是否沉默地產生錯檔；C 是否在 gate 明確失敗、是否指出漂移位置、
%      重新推導需要多少成本。這是 fail-loud 設計原則的直接證據。

\section{Ablations}
\label{sec:eval:ablation}

% TODO 拿掉 structural gates／只留單一 oracle（證明 §4.3 的互補性）／拿掉 repair loop／
%      換模型（至少兩個，證明結論不綁單一供應商）。

\section{Results}
\label{sec:eval:results}

% TODO 實驗尚未執行。呈現形式先定：
%      Table 主結果（三 arm x 全部指標）
%      Table 消融
%      Fig. break-even 曲線（三條線）
%      Fig. 漂移情境對照
%      失敗案例定性分析——每一類失敗各給一個具體例子，這部分最能說服領域讀者
%
% TODO 主表的形狀已定（2026-08-15，評測協定 §5.3）：**逐模板一列，沒有單一總平均**。
%      跨載體合併本來就被禁止，所以主表是「每個評測模板一列 × 三個 arm」。
%      不要在正文寫出一個把 S3/S4/S5 平均起來的數字，那個數字沒有意義。
% TODO 外部來源模板（SpreadsheetBench 衍生，若採用）**獨立一小節，不進主表平均**。
%      理由要寫出來：S3/S4/S5 之間的差異是主困難不同（能力軸）；外部來源模板回答的是
%      「結論會不會只是因為 benchmark 是自己設計的」（效度探針）。把效度證據平均進
%      能力結論，兩邊都變模糊。它的 n 小是刻意的——要的是存在性陳述，不是檢定。
%      預期它的數字會比自建模板低（真實檔案較髒），而分開報才能把那句話寫成
%      「這一列是壓力測試」而不是被動解釋「為什麼被平均掉」。

\subsection{Where the Randomness Went}
\label{sec:eval:randomness}

% TODO 誠實的一小節，緊接在可靠性數字之後：
%      本方法沒有消除隨機性，只是把它從 O(N) 次移到 1 次，
%      並且在那 1 次上加了 gate 與 oracle。
%      這個交換為什麼值得——以及在什麼條件下不值得（N 太小、derive 成功率太低）。
%      後者的展開在 §6.2。

\section{Case Study: Engineering Drawings}
\label{sec:eval:dwg}

% TODO DWG 獨立報告，不併入主結果表。
%      價值在於證明方法在表徵最不友善、且交付物有法規與生產責任的載體上一樣成立。
%      現有規模：1 份基準模板、1 份規則文件、5 個回歸案例、4 個標準答案圖、0 個漂移變體。

\section{Threats to Validity}
\label{sec:eval:threats}

Each item below states what is threatened, what was measured, and which arm the bias favours.
Items whose direction favours the proposed method are listed first, because those are the ones
a reader has no independent way to discover.

\subsection{Threats Biased Towards the Proposed Method}
\label{sec:eval:threats:favourable}

\paragraph{T1 --- development happened on evaluation templates.}
The template-level separation described in Section~\ref{sec:eval:split} does not hold: the two
development templates the protocol proposes were never built, and the prompts, gates, and
intermediate representation were tuned against the calibration instances of the evaluation
templates themselves. Three representation primitives exist specifically because the compiled
arm could not otherwise produce X1. The bias favours Arm C, and it is not quantified --- the
only way to quantify it is to build the missing development templates. The two bounding
measurements are given in Section~\ref{sec:eval:split}.

\paragraph{T2 --- the cell-datatype exemption.}
The compiler stringifies every value before writing it, so a numeric field lands in a
spreadsheet as text. The annotation rules therefore declare \texttt{presentation.datatype} as
\texttt{any} for every field. Without that exemption Arm C would lose a point on every numeric
field, and the study would be measuring a write-side type policy rather than the method. With
it, a real defect --- numbers stored as text --- is not scored, and a baseline that writes
correct types earns nothing for doing so. This is the third admissible response to a gap: the
claim is kept, labelled normative, and the direction disclosed. It favours Arm C.

\paragraph{T3 --- the rule tables are ours.}
Even where the template is a genuine third-party artifact, the machine-readable rule
applicability manifest that supplies the coverage denominator was written for this study. It
weakens the objection that the templates are self-designed; it does not weaken the objection
that the rules were written by someone who knew what the method can express. Any use of an
externally sourced template must state this in the same breath, or it reads as padding.

\paragraph{T4 --- three primitives rest on a single example.}
Artifact collections, stacked records, and equality merge between row sets were each validated
on one dataset, and it is the same dataset in all three cases. The key--value extent added for
the DOCX dataset has the same property: a survey of six real government forms found no second
instance of the shape. The failure mode this predicts is not a red test but a different form
that silently does not fit, and no measurement in this thesis would detect it.

\subsection{Threats Biased Against the Proposed Method}
\label{sec:eval:threats:unfavourable}

\paragraph{T5 --- package preservation was, until recently, unmeasurable for Arm C.}
The integrity channel treats drawings and printer settings as critical parts of a spreadsheet
package, which is why the government form was chosen as a dataset in the first place. Until it
was corrected, the production write-back for XLSX serialized the workbook through a
spreadsheet model, and that operation drops exactly those two parts --- the same two parts the
mutation suite's integrity negative control drops, because the negative control \emph{is} a
round trip through the same library. Arm C therefore scored 0 of 5 on integrity by
construction, and the channel could not distinguish a correct run from a deliberately
corrupted one.

The response was the first admissible one: the write-back was replaced by a package editor
that edits the one worksheet part that must change and copies every other byte through. The
measurement after that change, on the same five instances, is 5 of 5 with no preservation
violation, while re-serializing the same outputs through the spreadsheet model still yields 0
of 5. Narrowing the channel to fit the writer was considered and rejected, because it would
have been a reduction of the thesis claim to match the implementation. The threat is recorded
rather than deleted, because every integrity figure produced before that correction is an
artifact of the writer and not a result.

\paragraph{T6 --- the DOCX write path has not had the same audit.}
The correction in T5 covers XLSX only. The layout-bearing parts of a Word package have not
been enumerated part by part, and S4 is the only DOCX evaluation dataset. Any DOCX integrity
figure in this thesis therefore carries an unmeasured residual.

\paragraph{T7 --- an unpublishable specification is not necessarily a failed one.}
The publish gate blocks on any unresolved item, and the unresolved set has repeatedly been
found to contain entries that are not defects. Two counters existed for the same quantity ---
the agent's own self-reported list and the gate's canonical set --- and they were never
compared: on the best X1 specification they read 13 and 47. Thirty-four of the 47 were detail
cells that a repeat block writes correctly, admitted to the gate's input list because the
author-time oracle reads expected text from that same list for a different purpose. Any count
of unresolved items must be decomposed before it is interpreted, and a headline
``unpublishable'' figure without that decomposition would understate the method.

\subsection{Threats of Undetermined or Mixed Direction}
\label{sec:eval:threats:mixed}

\paragraph{T8 --- dataset scale.}
Two evaluation templates exist and a third is planned, against an original intention of three
to five per carrier. Conclusions must be descriptive; no statistical significance may be
claimed, and no single average may be taken across carriers.

\paragraph{T9 --- the inclusion criterion is itself skewed.}
The datasets admit only tasks whose target row count is fixed, because the write side
deliberately does not insert rows, and they exclude formula results. Both exclusions coincide
with boundaries of the proposed method, so the excluded region may contain tasks the baselines
can do and the compiled method cannot. This is the second admissible response --- an explicit
expressiveness boundary --- and it is disclosed rather than argued away.

\paragraph{T10 --- what the untouched channel does not reach.}
Untouched-region comparison covers cells and paragraphs. It does not reach footnotes,
endnotes, comments, or other independent package parts, and styles and formula results are
outside the first version of the scoring policy by design, to avoid confounding correctness
with whether a spreadsheet application refreshed its cached values.

\paragraph{T11 --- model dependence.}
Every derivation reported here used one model family. The ablation in
Section~\ref{sec:eval:ablation} is the instrument for this and has not been run; until it is,
no claim in this thesis is established as vendor-independent.

\paragraph{T12 --- the correctness of the reference outputs.}
Discussed in Section~\ref{sec:eval:selfvalidation}. Mutation testing bounds the evaluator, not
the generator; the residual is covered only by a probabilistic spot audit.

\paragraph{T13 --- two responses to the pilot that were rejected, and why.}
After the pilot in Section~\ref{sec:eval:rq} showed the per-instance baseline to be
semantically accurate, two adjustments were available and both were declined. Increasing task
difficulty until the baseline begins to fail would tune the benchmark towards producing the
desired comparison, and substituting a weaker model would amount to selecting the opponent.
Recording them here is part of the claim: the emphasis of this thesis moved to consistency and
cost structure \emph{because} the baseline was accurate, not because the tasks were made
harder.

\paragraph{Status: implementation-backed for T1, T2, T5, T6, T7, T9, T10, T12; normative for
T3, T4, T8, T11, T13.}
The items marked implementation-backed each name a measurement in this repository. The others
are stated positions or unrun experiments, and are labelled so that no reader takes an
intention for a result.

% =====================================================================
% 【剩下的寫作規格】2026-08-23
% 已寫成正文：§5.1／§5.2（含三個小節）／§5.3／§5.5／§5.10。
% 仍是 % TODO 的：§5.4 Metrics（含 5.4.1 rerun budget）、§5.6 Drift、§5.7 Ablations、
%   §5.8 Results、§5.9 DWG case study。各節的註記就是它們的寫作規格，不要另開一份。
%
% TODO §5.3 的 Status 段承諾「A/B 的實作、預算、素材清單要完整交代」——
%      live adapter 建好之後要補一小節寫預算（重試上限、wall-clock 上限的實際數值），
%      現在寫會是憑空的數字。
% TODO §5.10 的 T5 與 §4.7 的 D4／D5 互相引用：三處講的是同一組修法的三個面向，
%      改任何一處都要同步。已於 2026-08-23 對齊過一次。
% TODO §5.8 Results 動筆時，§5.1 的 pilot 數字（69.2%）不可升格成 Arm A 的結果，
%      它不是走 harness 跑的。§5.3 的 Status 段已經寫了這句，兩處不可打架。
% =====================================================================
